\documentclass[aspectratio=169,11pt]{beamer}

\usepackage[T1]{fontenc}
\usepackage[utf8]{inputenc}
\usepackage[brazil]{babel}
\usepackage{lmodern}
\usepackage{booktabs}
\usepackage{array}
\usepackage{xcolor}

\definecolor{ufjfBlue}{HTML}{12355B}
\definecolor{deepTeal}{HTML}{0B6E69}
\definecolor{warmGray}{HTML}{F4F1EC}
\definecolor{textDark}{HTML}{202124}
\definecolor{mutedText}{HTML}{5F6368}
\definecolor{accentGold}{HTML}{C99700}

\setbeamercolor{normal text}{fg=textDark,bg=warmGray}
\setbeamercolor{frametitle}{fg=white,bg=ufjfBlue}
\setbeamercolor{title}{fg=ufjfBlue}
\setbeamercolor{subtitle}{fg=deepTeal}
\setbeamercolor{author}{fg=textDark}
\setbeamercolor{institute}{fg=mutedText}
\setbeamercolor{date}{fg=mutedText}
\setbeamercolor{itemize item}{fg=deepTeal}
\setbeamercolor{itemize subitem}{fg=deepTeal}
\setbeamercolor{block title}{fg=white,bg=deepTeal}
\setbeamercolor{block body}{fg=textDark,bg=white}
\setbeamercolor{block title alerted}{fg=white,bg=accentGold}
\setbeamercolor{block body alerted}{fg=textDark,bg=white}
\setbeamercolor{alerted text}{fg=accentGold}

\setbeamerfont{normal text}{size=\footnotesize}
\setbeamerfont{itemize/enumerate body}{size=\footnotesize}
\setbeamerfont{itemize/enumerate subbody}{size=\footnotesize}
\setbeamerfont{block title}{size=\footnotesize,series=\bfseries}
\setbeamerfont{block body}{size=\footnotesize}
\setbeamerfont{block title alerted}{size=\footnotesize,series=\bfseries}
\setbeamerfont{block body alerted}{size=\footnotesize}

\setbeamertemplate{navigation symbols}{}
\setbeamertemplate{itemize item}{\raisebox{0.25ex}{\color{deepTeal}\rule{0.55ex}{0.55ex}}}
\setbeamertemplate{itemize subitem}{\raisebox{0.25ex}{\color{deepTeal}\rule{0.38ex}{0.38ex}}}
\setbeamertemplate{blocks}[rounded][shadow=false]
\setbeamertemplate{footline}{
  \leavevmode%
  \begin{beamercolorbox}[wd=\textwidth,ht=2.6ex,dp=1.1ex,leftskip=0.8em,rightskip=0.8em]{author in head/foot}
    \color{mutedText}\scriptsize Metodologia Científica -- Análise epistemológica evolucionária\hfill\insertframenumber/\inserttotalframenumber
  \end{beamercolorbox}
}

\setbeamertemplate{frametitle}{
  \nointerlineskip
  \begin{beamercolorbox}[wd=\textwidth,ht=4.2ex,dp=1.4ex,leftskip=0.9em]{frametitle}
    \usebeamerfont{frametitle}\insertframetitle
  \end{beamercolorbox}
}

\setbeamertemplate{title page}{
  \vspace*{0.55cm}
  {\color{accentGold}\rule{0.18\textwidth}{2.5pt}}
  \vspace{0.55cm}

  \begin{minipage}{0.90\textwidth}
    {\usebeamerfont{title}\Large\bfseries\inserttitle\par}
    \vspace{0.28cm}
    {\usebeamerfont{subtitle}\large\insertsubtitle\par}
    \vspace{0.80cm}
    {\large\insertauthor\par}
    \vspace{0.20cm}
    {\normalsize\insertinstitute\par}
    \vspace{0.55cm}
    {\normalsize\insertdate\par}
  \end{minipage}

  \vfill
  {\color{ufjfBlue}\rule{\textwidth}{2.5pt}}
}

\newcommand{\slidecite}[1]{%
  \vfill
  {\color{mutedText}\footnotesize\raggedright Fonte: #1\par}
}

\title{Análise Epistemológica}
\subtitle{\textit{Wi-Fi RTT Indoor Positioning}\\\textit{Using Visibility Matching With NLOS Receptions}}
\author{Programa de Pós-Graduação em Ciência da Computação -- UFJF}
\institute{Disciplina: Metodologia Científica -- Prof. Marcelo Bernardes Vieira}
\date{\today}

\begin{document}

\begin{frame}[plain]
  \titlepage
\end{frame}

\begin{frame}{Juízo central}
  \footnotesize
  \begin{itemize}
    \item Este é um artigo \textbf{evolucionário}: refina o paradigma vigente de posicionamento por rádio em ambientes internos, mas não inaugura um novo.
    \item Não cria uma nova ideia central; combina ideias já disponíveis.
    \item Tem genealogia rastreável: RADAR, Wi-Fi fingerprinting, RTT, NLOS e \textit{shadow/visibility matching}.
    \item O ganho quantitativo é forte, mas da ordem de refinamento técnico, não de ruptura.
    \item Isso não diminui seu valor científico: engenharia avança muito por ciência normal bem feita.
    \item Representa o estado da arte recente em posicionamento Wi-Fi RTT sob condições NLOS.
  \end{itemize}

  \slidecite{Kuhn; Lyu et al. (2025); Bahl e Padmanabhan (2000).}
\end{frame}

\begin{frame}{Motivação e relevância}
  \footnotesize
  \begin{itemize}
    \item Trabalho recente: publicado em junho de 2025.
    \item Publicado no \textit{IEEE Internet of Things Journal}, periódico Q1 de alto impacto.
    \item Dá continuidade a trabalhos do IPN-Lab (Hong Kong Polytechnic University).
    \item Guohao Zhang, com h-index 27 e mais de 2.300 citações, é o autor-âncora.
    \item Fronteira ativa: \textit{NLOS identification} + Wi-Fi RTT.
    \item Indicadores bibliométricos ainda iniciais: 5--6 citações, Impact Factor JCR'24 de 8,9 e 5-year Impact Factor de 9,6.
  \end{itemize}

  \vspace{0.18cm}
  \begin{block}{Leitura bibliométrica}
    As poucas citações refletem a curta janela de circulação. O ponto forte, por ora, é a posição em uma frente técnica recente, não uma influência já consolidada.
  \end{block}

  \slidecite{Lyu et al. (2025); \textit{IEEE Internet of Things Journal}; Google Scholar.}
\end{frame}

\begin{frame}{Problema e situação de contorno}
  \footnotesize
  \vspace{-0.18cm}
  \begin{block}{Pergunta que organiza o artigo}
    Como posicionar um usuário em ambiente interno por Wi-Fi RTT quando há NLOS severo, poucos APs e medições com outliers?
  \end{block}

  \vspace{0.24cm}
  \setlength{\tabcolsep}{3pt}
  \renewcommand{\arraystretch}{0.92}
  \begin{tabular}{>{\bfseries\raggedright\arraybackslash}p{0.22\textwidth}>{\raggedright\arraybackslash}p{0.68\textwidth}}
    \toprule
    Elemento & Recorte epistemológico \\
    \midrule
    Objeto & Dispositivo, APs Wi-Fi, RTT/RSSI e planta digital. \\
    Fenômeno & Propagação LOS/NLOS, atenuação, multipercurso e erro de alcance. \\
    Domínio & Ambientes internos com infraestrutura Wi-Fi e planta conhecida. \\
    Situação de contorno & NLOS severo, poucos APs, posições conhecidas e planta digital disponível. \\
    Condições de operação & RTT/RSSI, candidatos em grade e visibilidade por interseção linha--parede. \\
    Métrica & Erro de posição, média, mediana, desvio padrão e CDF. \\
    \bottomrule
  \end{tabular}

\end{frame}

\begin{frame}{O deslocamento epistemológico}
  \footnotesize
  \begin{columns}[T]
    \begin{column}{0.47\textwidth}
      \textbf{Antes: NLOS como erro}
      \begin{itemize}
        \item O NLOS degrada o RTT e gera erro de alcance.
        \item O sinal bloqueado ou refletido aparece como perturbação.
        \item A tarefa é mitigar, corrigir ou atenuar essa perturbação.
      \end{itemize}
    \end{column}
    \begin{column}{0.49\textwidth}
      \textbf{Neste artigo: NLOS como informação}
      \begin{itemize}
        \item O NLOS ajuda a inferir se um AP está visível ou invisível.
        \item A visibilidade dos APs torna-se uma característica saliente.
        \item A planta digital transforma essa característica em restrição espacial.
      \end{itemize}
    \end{column}
  \end{columns}

  \vspace{0.25cm}
  \begin{alertblock}{Mudança conceitual}
    O NLOS deixa de ser apenas erro de medição e passa a revelar uma característica saliente: a visibilidade dos APs.
  \end{alertblock}

\end{frame}

\begin{frame}{Hipótese e ensaio de solução}
  \footnotesize
  \begin{itemize}
    \setlength{\itemsep}{0.25cm}
    \item \textbf{Hipótese física}: paredes alteram RTT e RSSI de modo informativo.
    \item \textbf{Hipótese geométrica}: a planta digital permite prever visibilidade dos APs.
    \item \textbf{Hipótese estatística}: RTT/RSSI permitem inferir visibilidade observada.
    \item \textbf{Hipótese computacional}: a melhor posição é a que maximiza a coerência entre alcance e visibilidade.
    \item \textbf{Ensaio de solução}: modelo multiwall de RTT + visibilidade simulada + visibilidade observada + busca em grade.
  \end{itemize}

\end{frame}

\begin{frame}{Conteúdo lógico e empírico}
  \footnotesize
  \begin{columns}[T]
    \begin{column}{0.47\textwidth}
      \textbf{Conteúdo lógico}
      \begin{itemize}
        \item Equações de alcance.
        \item Modelo multiwall.
        \item Interseção linha--parede.
        \item XNOR de visibilidade.
        \item Função de pontuação.
      \end{itemize}
    \end{column}
    \begin{column}{0.49\textwidth}
      \textbf{Conteúdo empírico}
      \begin{itemize}
        \item Três ambientes reais na PolyU.
        \item Três APs Google Nest Wi-Fi.
        \item RTT/RSSI a 3 Hz.
        \item Ground truth por LiDAR.
        \item Comparação com baselines.
      \end{itemize}
    \end{column}
  \end{columns}

  \vspace{0.25cm}
  \begin{alertblock}{Papel no artigo}
    O modelo formal define o que o método permite inferir; os experimentos indicam quais inferências se sustentam como observação nos cenários testados.
  \end{alertblock}

  \slidecite{Lyu et al. (2025), Seção V.}
\end{frame}

\begin{frame}{Evidência experimental}
  \footnotesize
  \setlength{\tabcolsep}{2pt}
  \renewcommand{\arraystretch}{0.95}
  \begin{tabular}{lrrrrrr}
    \toprule
    Método & Corr. média & Corr. mediana & Lounge média & Lounge mediana & Suíte média & Suíte mediana \\
    \midrule
    TRI & 1,399 & 1,187 & 2,415 & 2,683 & 1,672 & 1,688 \\
    PF & 1,462 & 1,314 & 2,444 & 2,588 & 1,854 & 1,866 \\
    SMACOF & 1,402 & 1,166 & 1,601 & 1,473 & 1,814 & 1,772 \\
    \textbf{VM} & \textbf{0,681} & \textbf{0,673} & \textbf{0,801} & \textbf{0,785} & \textbf{0,695} & \textbf{0,823} \\
    \bottomrule
  \end{tabular}

  \vspace{0.30cm}
  \begin{block}{Corroboração, não prova}
    A evidência sustenta a vantagem do VM nos cenários testados, mas não autoriza concluir superioridade universal.
  \end{block}

  \slidecite{Lyu et al. (2025), Tabela I.}
\end{frame}

\begin{frame}{Falseabilidade e reprodutibilidade}
  \footnotesize
  \begin{columns}[T]
    \begin{column}{0.47\textwidth}
      \textbf{Como falsear}
      \begin{itemize}
        \item Testar layouts e materiais diferentes.
        \item Alterar quantidade e distribuição de APs.
        \item Usar plantas incompletas ou desatualizadas.
        \item Comparar com métodos posteriores.
      \end{itemize}
    \end{column}
    \begin{column}{0.49\textwidth}
      \textbf{Como reproduzir}
      \begin{itemize}
        \item Há equipamentos, frequência e métricas.
        \item Há equações e baselines.
        \item Faltam código, dataset e parâmetros completos.
        \item A réplica exata depende da planta e do ambiente.
      \end{itemize}
    \end{column}
  \end{columns}

  \vspace{0.25cm}
  \begin{alertblock}{Leitura popperiana}
    O artigo oferece uma teoria operacional criticável: boa ciência, mas ainda dependente de condições auxiliares.
  \end{alertblock}

\end{frame}

\begin{frame}{Análise crítica}
  \footnotesize
  \begin{tabular}{>{\bfseries\raggedright\arraybackslash}p{0.24\textwidth}>{\raggedright\arraybackslash}p{0.66\textwidth}}
    \toprule
    Dimensão & Avaliação crítica \\
    \midrule
    Problema & Bem delimitado: NLOS severo, poucos APs e outliers de distância. \\
    Hipótese & Plausível: NLOS deixa de ser apenas erro e passa a carregar informação espacial. \\
    Método & Coerente, mas dependente de planta digital correta e ajuste de parâmetros. \\
    Evidência & Boa nos três cenários testados, com comparação com baselines. \\
    Generalização & Limitada para outros edifícios, materiais, densidades de APs e mudanças ambientais. \\
    Reprodutibilidade & Parcial: há equações e métricas, mas faltam código, dataset e parâmetros completos. \\
    \bottomrule
  \end{tabular}

  \vspace{0.15cm}
  \begin{block}{Crítica central}
    O método funciona bem nos cenários testados, mas ainda não há evidência de robustez em ambientes internos variados.
  \end{block}
\end{frame}

\begin{frame}{Limites e alcance da contribuição}
  \footnotesize
  \begin{itemize}
    \item \textbf{Força epistemológica}: muda o estatuto do NLOS, de perturbação experimental para pista informacional.
    \item \textbf{Força metodológica}: articula modelo físico, geometria da planta, inferência LOS/NLOS e validação empírica.
    \item \textbf{Limite de representação}: o ambiente real é reduzido a planta digital, APs, paredes e candidatos em grade.
    \item \textbf{Limite operacional}: portas, móveis, pessoas, materiais e erros na planta podem alterar a visibilidade esperada.
    \item \textbf{Impacto}: por ser recente, o baixo número de citações não é fraqueza científica; apenas impede afirmar uma influência já consolidada.
  \end{itemize}
\end{frame}

\begin{frame}{Síntese final}
  \footnotesize
  \begin{itemize}
    \item O artigo avança por aproximação sucessiva, não por ruptura.
    \item O centro epistemológico é a mudança de estatuto do NLOS.
    \item Opera predominantemente em chave Kantiana: sintetiza o modelo formal Leibniziano/Multiwall com dados empíricos Lockeanos de RSSI/RTT.
    \item A evidência é forte para os cenários testados e limitada para generalização.
    \item A crítica principal recai sobre o alcance externo, a dependência da planta, a calibração e os artefatos de reprodução.
  \end{itemize}

  \vspace{0.28cm}
  \begin{alertblock}{Frase de fechamento}
    O artigo não muda a pergunta da área; ele melhora a resposta ao transformar uma fraqueza experimental em uma pista epistemologicamente útil.
  \end{alertblock}

\end{frame}

\end{document}
